\ExerciseSection

\begin{enumerate}
\def\labelenumi{\arabic{enumi})}
\tightlist
\item
  \begin{enumerate}
  \def\labelenumii{\alph{enumii})}
  \tightlist
  \item
    Every zero-dimensional space is completely regular. The non-normal locally compact spaces defined in § 4, Exercise 26 b) and § 5, Exercise 15 are zero-dimensional.
  \end{enumerate}
\end{enumerate}

\begin{enumerate}
\def\labelenumi{\alph{enumi})}
\setcounter{enumi}{1}
\item
  A topological space X is said to be strongly zero-dimensional if for each closed subset A of X and each neighbourhood U of A, there is an open and closed neighbourhood of A contained in U. Every strongly zero-dimensional space is normal. A normal space is strongly zero-dimensional if and only if its Stone-Čech compactification (§ 1, Exercise 7) is totally disconnected {[}use Exercise 17 c) of § 4{]}.
\item
  Let X be a strongly zero-dimensional Hausdorff space and let \((\mathbf{U}_{n})\) be an increasing sequence of non-empty open sets in X. Show that there
\end{enumerate}

\[
\left(\mathbf {G} _ {n}\right)
\]

\[
\mathbf {x},
\]

\[
\mathbf {G} _ {n} \subset \mathbf {U} _ {n}
\]

\[
n,
\]

\[
\bigcup_ {n} \mathrm{G} _ {n} = \bigcup_ {n} \mathrm{U} _ {n}
\]

(define the \(G_{n}\) by induction). If, moreover, X is perfectly normal (§ 4, Exercise 8), every open set in X is the union of a countable family of mutually disjoint sets which are both open and closed in X.

\begin{enumerate}
\def\labelenumi{\alph{enumi})}
\setcounter{enumi}{3}
\tightlist
\item
  Let X be a normal space which is the union of a sequence \((\mathbf{A}_{n})\) of strongly zero-dimensional subspaces. Show that X is strongly zero-dimensional. {[}Let B, C be two disjoint closed subsets of X; define by induction two increasing sequences \((\mathbf{G}_{n}), (\mathbf{H}_{n})\) of open sets in X such that \(\overline{G}_{n} \cap \overline{H}_{n} \neq \varnothing\) , \(A_{n} \subset G_{n} \cup H_{n}\) , \(B \cap A_{n} \subset G_{n}\) , \(C \cap A_{n} \subset H_{n}\) .{]}
\end{enumerate}

¶ 2) a) Let X be a metrizable space. Show that the following properties are equivalent:

\(\alpha)\) There exists a metric \(d\) compatible with the topology of \(\mathbf{X}\) , with respect to which \(\mathbf{X}\) is an ultrametric space (§ 2, Exercise 4).

\(\beta)\) X is strongly zero-dimensional.

γ) There exists a sequence \((\mathfrak{B}_{n})\) of locally finite families of open and closed subsets of X such that \(B = \bigcup_{n} B_{n}\) is a base for the topology of X.

{[}To show that \(\alpha\) ) implies \(\beta\) ), note that if F is a closed set in an ultrametric space X, then the set of all \(x \in X\) such that \(d(x, F) = \rho > 0\) is both open and closed in X. To show that \(\beta\) ) implies \(\gamma\) ), use Lemma 3 of § 4, no. 5 and Exercise 1 c) of § 6. Finally, to establish that \(\gamma\) implies \(\alpha\) ), reduce to the case where \(\mathfrak{B}_n = (\mathrm{U}_{n,\lambda})_{\lambda \in L}\) , and if \(f_{n,\lambda}\) is the characteristic function of \(\mathrm{U}_{n,\lambda}\) , consider the mapping \(x \to f(x) = f_{n,\lambda}(x)\) of X into the product group \(\mathbf{G}^{\mathbf{N} \times \mathbf{L}}\) , where \(\mathbf{G} = \mathbf{Z}/2\mathbf{Z}\) (identified with the set \(\{0, 1\}\) ); for each \(n \in \mathbb{N}\) , let \(\mathbf{G}_n\) be the subgroup of G consisting of the elements whose \((k, \lambda)\) coordinates are zero for all \(\lambda \in L\) and all \(k < n\) ; show that if G is endowed with the group topology for which \((\mathrm{G}_n)\) is a fundamental system of neighbourhoods of o, then \(f\) is a homeomorphism of X onto a subspace of G; and complete the proof by noting that the topology of G can be defined by an invariant metric with respect to which G is an ultrametric space.{]}

\begin{enumerate}
\def\labelenumi{\alph{enumi})}
\setcounter{enumi}{1}
\tightlist
\item
  Deduce from a) that every zero-dimensional Hausdorff space X which has a countable base is a strongly zero-dimensional metrizable space. (Show by using § 2, no. 8, Proposition 13 that there exists a countable base formed of sets which are both open and closed.) Furthermore, X is then homeomorphic to a subspace of Cantor's triadic set K (cf.~§ 2, Exercise 12 (*).)
\end{enumerate}

¶ 3) a) Show that every totally disconnected subspace of R is zero-dimensional.

\begin{enumerate}
\def\labelenumi{\alph{enumi})}
\setcounter{enumi}{1}
\tightlist
\item
  Let K be Cantor's triadic set; every \(x \in K\) can be written uniquely in the form \(x = \sum_{k} 2/3^{n_k}\) , where \((n_k)_{k \geqslant 1}\) is a strictly increasing (finite or infinite) sequence of integers \textgreater0 (Chapter IV, § 8, Exercise 9); and define \(f(x)\) to be the number \(\sum_{k} (-1)^{n_k}/2^k\) . Let G be the graph of the function f in \(K \times R\) . Show that the metrizable space G is totally disconnected but is not zero-dimensional. {[}Note that the intersection of \(\overline{G}\) and the line \(\{0\} \times R\) contains an open interval containing the point \((0, 0)\) .{]}
\end{enumerate}

\par\noindent\textbf{4) Let \(\mathbf{X}\) be a metrizable space.}\par

\begin{enumerate}
\def\labelenumi{\alph{enumi})}
\item
  Show that the set of Borel subsets of X is the smallest subset F of \(\mathfrak{P}(X)\) which contains all closed subsets of X and is such that countable unions and countable intersections of sets of F belong to F. (Consider the subset G of F consisting of all \(A \in F\) such that \(X - A \in F\) .)
\item
  Show that the set of Borel subsets of X is the smallest subset F of \(\mathfrak{P}(X)\) which contains all open subsets of X and is such that every countable intersection of sets of F belongs to F, and every union of a sequence of mutually disjoint sets of F belongs to F (same method).
\item
  Every ordinal \(\alpha\) can be written uniquely in the form \(\alpha = \omega\beta + n\) , where \(n < \omega\) (Set Theory, Chapter III, § 2, Exercise 15); \(\alpha\) is said to be even or odd according as \(n\) is even or odd. For every countable ordinal \(\alpha\) , we define by transfinite induction sets of subsets \(\mathfrak{F}_{\alpha}(\mathbf{X})\) , \(\mathfrak{G}_{\alpha}(\mathbf{X})\) (or simply \(\mathfrak{F}_{\alpha}\) , \(\mathfrak{G}_{\alpha}\) ) as follows: \(\mathfrak{F}_0\) (resp. \(\mathfrak{G}_0\) ) is the set of all closed (resp. open) subsets of X; if \(\alpha\) is even and \(>0\) , \(\mathfrak{F}_{\alpha}\) (resp. \(\mathfrak{G}_{\alpha}\) ) is the set of all subsets of X which are countable intersections (resp. countable unions) of sets belonging to \(\bigcup_{\xi < \alpha} \mathfrak{F}_{\xi}\) (resp. \(\bigcup_{\xi < \alpha} \mathfrak{G}_{\xi}\) ); if \(\alpha\) is odd, then \(\mathfrak{F}_{\alpha}\) (resp. \(\mathfrak{G}_{\alpha}\) ) is the set of all subsets of X which are countable unions (resp. countable intersections) of sets belonging to \(\bigcup_{\xi < \alpha} \mathfrak{F}_{\xi}\) (resp. \(\bigcup_{\xi < \alpha} \mathfrak{G}_{\xi}\) ). Show that the union of the \(\mathfrak{F}_{\alpha}\) (resp. \(\mathfrak{G}_{\alpha}\) ) is the set of Borel subsets of X {[}use \(a\) {]}. Deduce that if X is of countable type, the set of Borel subsets of X has power at most equal to the power of the continuum.
\end{enumerate}

\(d\) ) The relation \(\mathbf{A}\in\mathfrak{F}_{\alpha}\) is equivalent to \(\mathbb{C}\mathbf{A}\in\mathfrak{G}_{\alpha}\) . Every finite union (resp. intersection) of sets of \(\mathfrak{F}_{\alpha}\) (resp. \(\mathfrak{G}_{\alpha}\) ) belongs to \(\mathfrak{F}_{\alpha}\) (resp. \(\mathfrak{G}_{\alpha}\) ). We have \(\mathfrak{F}_{\alpha}\subset\mathfrak{F}_{\alpha+1}\cap\mathfrak{G}_{\alpha+1}\) and \(\mathfrak{G}_{\alpha}\subset\mathfrak{F}_{\alpha+1}\cap\mathfrak{G}_{\alpha+1}\) (use transfinite induction).

\begin{enumerate}
\def\labelenumi{\alph{enumi})}
\setcounter{enumi}{4}
\tightlist
\item
  Let \(f\) be a continuous mapping of \(\mathbf{X}\) into a metrizable space \(\mathbf{Y}\) . Show that \(f^{-1}(\mathfrak{F}_{\alpha}(\mathbf{Y})) \subset \mathfrak{F}_{\alpha}(\mathbf{X})\) and \(f^{-1}(\mathfrak{G}_{\alpha}(\mathbf{Y})) \subset \mathfrak{G}_{\alpha}(\mathbf{X})\) .
\end{enumerate}

\(f)\) If \(\mathbf{X}\) and \(\mathbf{Y}\) are two metrizable spaces and \(\mathbf{A} \in \mathfrak{F}_{\alpha}(\mathbf{X}), \mathbf{B} \in \mathfrak{F}_{\alpha}(\mathbf{Y})\) , {[}resp. \(\mathbf{A} \in \mathfrak{G}_{\alpha}(\mathbf{X}), \mathbf{B} \in \mathfrak{G}_{\alpha}(\mathbf{Y})\) {]}, then \(\mathbf{A} \times \mathbf{B} \in \mathfrak{F}_{\alpha}(\mathbf{X} \times \mathbf{Y})\) {[}resp. \(\mathbf{A} \times \mathbf{B} \in \mathfrak{G}_{\alpha}(\mathbf{X} \times \mathbf{Y})\) {]}.

\begin{enumerate}
\def\labelenumi{\alph{enumi})}
\setcounter{enumi}{6}
\item
  Let \((\mathbf{X}_n)\) be a sequence of metrizable spaces. Show that if \(\mathbf{A}_n\) is a Borel set in \(\mathbf{X}_n\) for each \(n\) , then \(\prod_{n} \mathbf{A}_n\) is a Borel set in \(\prod_{n} \mathbf{X}_n\) .
\item
  Let \(\alpha > 0\) be an even (resp. odd) countable ordinal and let \((\mathbf{A}_n)\) be a countable covering of \(\mathbf{X}\) formed of sets belonging to \(\mathfrak{G}_{\alpha}\) (resp. \(\mathfrak{F}_{\alpha}\) ). Show that there is a partition \((\mathbf{B}_n)\) of \(\mathbf{X}\) such that \(\mathbf{B}_n \subset \mathbf{A}_n\) and \(\mathbf{B}_n \in \mathfrak{F}_{\alpha} \cap \mathfrak{G}_{\alpha}\) for each \(n\) {[}use \(d\) {]}.
\end{enumerate}

¶ 5) a) Let J be the Polish space \(\mathbf{N}^{\mathbf{N}}\) (N carrying the discrete topology) and let X be a metrizable space of countable type. Show that for each countable ordinal \(\alpha\) there is a subset \(\mathbf{G}_{\alpha}\) of \(\mathbf{J} \times \mathbf{X}\) such that (i) \(\mathbf{G}_{\alpha} \in \mathfrak{G}_{\alpha}\) ( \(\mathbf{J} \times \mathbf{X}\) ) (Exercise 4) and (ii) for every subset \(\mathbf{U} \in \mathfrak{G}_{\alpha}(\mathbf{X})\) there exists \(z \in \mathbf{J}\) such that \(\mathbf{G}_{\alpha}(z) = \mathbf{U}\) . {[}We may proceed as follows by transfinite induction: the spaces J and \(\mathbf{J}^{\mathbf{N}}\) being homeomorphic, let \(f\) be a homeomorphism of J onto \(\mathbf{J}^{\mathbf{N}}\) , and for each integer \(n\) let \(f_{n} = \mathrm{pr}_{n} \circ f\) . If \((\mathbf{A}_{n})\) is a countable base of X containing the empty set, take \(\mathbf{G}_{0}\) such that \(\mathbf{G}_{0}(z) = \bigcup_{n} \mathbf{A}_{f_{n}(z)}\) for each \(z \in \mathbf{J}\) . For each countable ordinal \(\alpha\) , take \(\mathbf{G}_{\alpha+1}\) to be such that \(\mathbf{G}_{\alpha+1}(z) = \bigcap_{n} \mathbf{G}_{\alpha}(f_{n}(z))\) if \(\alpha\) is even, and \(\mathbf{G}_{\alpha+1}(z) = \bigcup_{n} \mathbf{G}_{\alpha}(f_{n}(z))\) if \(\alpha\) is odd. Finally, if \(\alpha\) has no predecessor and if \((\lambda_{n})\) is an increasing sequence of ordinals such that \(\alpha = \sup_{n} \lambda_{n}\) , take \(\mathbf{G}_{\alpha}\) to be such that \(\mathbf{G}_{\alpha}(z) = \bigcup_{n} \mathbf{G}_{\lambda_{n}}(f_{n}(z))\) . Use the fact that \(f_{n}\) is continuous on J.{]}

\begin{enumerate}
\def\labelenumi{\alph{enumi})}
\setcounter{enumi}{1}
\tightlist
\item
  In particular take \(\mathbf{X} = \mathbf{J}\) . Show that the set \(\mathrm{pr}_1(\mathbf{G}_\alpha \cap \Delta)\) , where \(\Delta\) is the diagonal of \(\mathbf{J} \times \mathbf{J}\) , belongs to \(\mathfrak{G}_{\alpha}(\mathbf{J})\) but not to \(\mathfrak{F}_{\alpha}(\mathbf{J})\) {[}argue by contradiction, using \(a\) {]}.
\end{enumerate}

\begin{enumerate}
\def\labelenumi{\arabic{enumi})}
\setcounter{enumi}{5}
\tightlist
\item
  \begin{enumerate}
  \def\labelenumii{\alph{enumii})}
  \tightlist
  \item
    Let J be the Polish space \(N^{N}\) . Show that if X is any Souslin space, there is a continuous surjection of J onto X {[}reduce to the case where X is Polish and consider a sifting of X by a sieve ( \(C_{n}\) ) where all the \(C_{n}\) are infinite{]}.
  \end{enumerate}
\end{enumerate}

\begin{enumerate}
\def\labelenumi{\alph{enumi})}
\setcounter{enumi}{1}
\item
  Let X be a metrizable space. Show that if A is any Souslin subspace of X, there is a closed subset F of \(J \times X\) such that \(\mathbf{A} = \operatorname{pr}_{2}(\mathbf{F})\) {[}use a){]}.
\item
  Let X be a metrizable space of countable type; let \(L = J \times X\) and let F be a closed subset of \(J \times L\) such that for each closed subset
\end{enumerate}

M of L there exists \(z \in J\) such that \(F(z) = M\) {[}Exercise 5 a){]}. Show that, for each Souslin subspace S of X, there exists \(z \in J\) such that \(\operatorname{pr}_{2}(\mathbf{F}(z)) = \mathbf{S}\) . Deduce that, if X = J, the set T of all \(z \in J\) such that \(z \in \operatorname{pr}_{2}(\mathbf{F}(z))\) is a Souslin set, but that J --- T is not a Souslin set {[}same reasoning as in Exercise 5 b){]} (*).

\begin{enumerate}
\def\labelenumi{\arabic{enumi})}
\setcounter{enumi}{6}
\tightlist
\item
  \begin{enumerate}
  \def\labelenumii{\alph{enumii})}
  \tightlist
  \item
    Show that every zero-dimensional Polish space is homeomorphic to a closed subspace of the product \(J = N^{N}\) (consider a strict sifting of X by sets which are both open and closed).
  \end{enumerate}
\end{enumerate}

\begin{enumerate}
\def\labelenumi{\alph{enumi})}
\setcounter{enumi}{1}
\item
  Let X be a zero-dimensional Polish space and let Y be a dense subspace of X which has no interior point and which is a countable intersection of open sets in X. Show that Y is homeomorphic to J. (Note that, in a metrizable zero-dimensional space, an open but not closed set is the union of an infinite countable sequence of pairwise disjoint sets which are both open and closed; use Theorem 1 of no. 1.) Deduce that every dense subspace of R, which has no interior point and is a countable intersection of open sets, is homeomorphic to J (note that such a set is contained in the complement of a countable dense subset D of R, and that CD is zero-dimensional).
\item
  If X is any uncountable zero-dimensional Polish space, there exists a partition of X into a countable set and a subspace homeomorphic to J {[}use b) and Exercise 11 of § 5{]}.
\end{enumerate}

\begin{enumerate}
\def\labelenumi{\arabic{enumi})}
\setcounter{enumi}{7}
\tightlist
\item
  Let X be a Souslin space and let f be a continuous mapping of X into a Hausdorff space Y. Show that if \(f(\mathbf{X})\) is uncountable, there exists a subspace A of X homeomorphic to Cantor's triadic set K (Chapter IV, § 2, no. 5) such that \(f|A\) is injective. {[}Reduce to the case where X is a complete metric space: show that there exists a sieve \(\mathbf{C} = (\mathbf{C}_{n}, p_{n})\) such that, for each n, \(C_{n}\) has \(2^{n}\) elements, and that there exists for each n a mapping \(\varphi_{n}\) of \(C_{n}\) into the set of all non-empty closed subsets of diameter \(\leqslant 2^{-n}\) in X, such that (i) \(\varphi_{n+1}(c) \subset \varphi_{n}(p_{n}(c))\) for all \(c \in C_{n+1}\) ; (ii) whenever c and \(c'\) are distinct elements of \(C_{n}\) , \(f(\varphi_{n}(c))\) and \(f(\varphi_{n}(c'))\) are disjoint. Use § 5, Exercise 11.{]} In particular, every uncountable Souslin space contains a subspace homeomorphic to K and therefore has the power of the continuum.
\end{enumerate}

¶ 9) a) Let X be a Polish space, and let R be an equivalence relation on X. Show that there exists a subset of X which is a countable intersection of open sets and which meets each equivalence class in exactly one point, provided one or the other of the following two hypotheses is satisfied :

\(\alpha)\) R is closed;

\(\beta)\) \(\mathbf{X}\) is locally compact and the graph of \(\mathbf{R}\) is closed in \(\mathbf{X} \times \mathbf{X}\) .

{[}For \(\alpha\) ), follow the method of the proof of Theorem 4 of no. 8, making use of the following remark: If \((\mathbf{G}_n)\) is a sequence of sets, each of which is a countable intersection of open sets, and if for each \(n\) there exists a neighbourhood \(\mathbf{V}_n\) of \(\mathbf{G}_n\) such that the family \((\mathbf{V}_n)\) is locally finite, then \(\bigcup_{n} \mathbf{G}_n\) is a countable intersection of open sets. Same method for \(\beta\) );

observe that the saturation of a compact set is now closed (cf.~Chapter I, § 10, Exercise 16).{]}

\begin{enumerate}
\def\labelenumi{\alph{enumi})}
\setcounter{enumi}{1}
\item
  Let X be a Polish space, let R be an equivalence relation on X which is both open and closed, and let \(\varphi: \mathbf{X} \to \mathbf{X} / \mathbf{R}\) be the canonical mapping. Show that there exists a subset A of X which is a countable intersection of open sets in X and is such that the restriction of \(\varphi\) to A is a homeomorphism of A onto \(\varphi(A)\) and such that \(\varphi(A)\) is dense in X/R and is a countable intersection of open sets. {[}With the notation of the proof of Theorem 4 of no. 8, show that we may assume that for each \(n\) the \(\varphi_n(c)\) have been defined in such a way that the image under \(\varphi\) of the union of the interiors of the \(h_n(c)\) , as \(c\) runs through \(C_n\) , is dense in X/R; use \(a\) ) and Theorem 1 of no. 1.{]}
\item
  Show that the conclusion of b) remains valid when X is locally compact and has a countable base and R is a closed equivalence relation on X all of whose equivalence classes are compact. {[}Reduce to case b) by the use of § 5, Exercise 26.{]}
\end{enumerate}

¶ 10) a) Let X be a metrizable space and let S be a Souslin subspace of X. Show that S is almost open in X (§ 5, Exercise 6). {[}Let P be a Polish space, g a continuous mapping of P onto S, (Cn, pn, φn) a sifting of P, and f the corresponding continuous mapping of L(C) onto P (in the notation of no. 5); let h = g∘f, and for each c ∈ Cn let qn(c) denote the subspace of L(C) consisting of sequences (ck) such that cn = c.~Let Fn(c) = h(qn(c)) and let Un(c) = g(φn(c)) ∃ Fn(c) for each c ∈ Cn; let Wn(c) denote the union of D(Un(c)) (§ 5, Exercise 3) and a meagre set in X, such that Fn(c) ⊂ Wn(c) ⊂ U̅n(c), so that Wn(c) is almost open in X. Show that

\[
\mathrm{V} _ {n} (c) = \mathrm{W} _ {n} (c) \cap \mathbb {C} \left(\bigcup_ {p _ {n} \left(c ^ {\prime}\right) = c} \mathrm{W} _ {n + 1} \left(c ^ {\prime}\right)\right)
\]

is meagre in X by noting that \(\mathbf{F}_n(c) = \bigcup_{p_n(c') = c}\mathbf{F}_{n + 1}(c')\) ; show also that the set \(\mathrm{W}_{n}(c)-\mathrm{F}_{n}(c)\) is contained in the union of the sets \(\mathrm{V}_{m}(d)\) where \(m\geqslant n\) and, for each m,d runs through the set of elements of \(C_{m}\) such that \(p_{nm}(d)=c\) . Hence show that \(\mathrm{F}_{n}(c)\) is almost open.{]}

\begin{enumerate}
\def\labelenumi{\alph{enumi})}
\setcounter{enumi}{1}
\tightlist
\item
  Give an example of a continuous mapping \(f\) of {[}0, 1{]} into itself and an almost open set \(\mathbf{B} \subset \mathbf{I}\) such that \(f(\mathbf{B})\) is not almost open. {[}Use Exercise 16 b) of Chapter IV, § 8, and show that we may take \(\mathbf{B}\) to be a subset of the Cantor set \(\mathbf{K}\) such that \(f(\mathbf{B}) = \mathbf{Z}\) has the property described in § 5, Exercise 18 d).{]}
\end{enumerate}

¶ 11) Let \(\mathbf{X} = \mathbf{R}^n\) and let \(\mathbf{M}\) be a closed subset of \(\mathbf{X}\) . A point \(x \in \mathbf{M}\) is said to be linearly accessible if there exists a point \(y \in \mathbf{X} - \mathbf{M}\) such that the open segment with extremities \(x, y\) is contained in \(\mathbf{X} - \mathbf{M}\) . Show that the set \(\mathbf{L} \subset \mathbf{M}\) of linearly accessible points of \(\mathbf{M}\) is a Souslin set. {[}Let \(d(x, y)\) be the Euclidean metric in \(\mathbf{X}\) , and for each \(z \in \mathbf{X}\) let \(f_z(x, y) = d(x, z) + d(z, y) - d(x, y)\) . Note that the set of all \((x, y, z) \in \mathbf{X} \times \mathbf{X} \times \mathbf{X}\) such that \(x \in \mathbf{M}, y \in \mathbf{X} - \mathbf{M}, z \in \mathbf{M}, z \neq x\) and \(f_z(x, y) = 0\) is a countable union of compact sets, and consequently so is its projection on the product of the first two factors.{]}

\begin{enumerate}
\def\labelenumi{\arabic{enumi})}
\setcounter{enumi}{11}
\tightlist
\item
  Let X be a metrizable space, and let f be a mapping of the set \(\mathcal{R}(\mathbf{X})\) of compact subsets of X into \(\overline{R}\) which satisfies the conditions \((\mathrm{CA}_{\mathrm{I}})\) (for two compact subsets of X) and \((\mathrm{CA}_{\mathrm{III}})\) . For each subset A of X, let \(f_{*}(\mathrm{A})\) denote the least upper bound of the numbers \(f(\mathrm{K})\) for all compact subsets K of A, and let \(f^{*}(\mathrm{A})\) denote the greatest lower bound of the \(f_{*}(\mathrm{U})\) for all open sets U containing A. A is said to be admissible with respect to f if \(f_{*}(\mathrm{A}) = f_{*}(\mathrm{A})\) .
\end{enumerate}

\begin{enumerate}
\def\labelenumi{\alph{enumi})}
\item
  Show that for every compact subset L of X and every \(\varepsilon > 0\) , there exists an open set \(U \supset L\) such that, for every compact subset K of X for which \(L \subset K \subset U\) , we have \(f(K) \leqslant f(L) + \varepsilon\) . (Argue by contradiction.) Deduce that every compact subset and every open subset of X is admissible with respect to f.
\item
  Suppose that \(f\) satisfies the following condition:
\end{enumerate}

(AL) For each pair of compact subsets K, K' of X, we have

\[
f (\mathbf {K} \cup \mathbf {K} ^ {\prime}) + f (\mathbf {K} \cap \mathbf {K} ^ {\prime}) \leqslant f (\mathbf {K}) + f (\mathbf {K} ^ {\prime}).
\]

Show that if \((\mathbf{K}_i)\) is a finite family of compact sets with the property that

\[
f \left(\bigcup_ {i} \mathrm{K} _ {i}\right) <   + \infty ,
\]

and if \((\mathbf{K}_i')\) is another finite family of compact sets indexed by the same set, such that \(\mathbf{K}_i' \subset \mathbf{K}_i\) and \(f(\mathbf{K}_i') > -\infty\) for each \(i\) , then we have

\[
f \left(\bigcup_ {i} \mathrm{K} _ {i}\right) - f \left(\bigcup_ {i} \mathrm{K} _ {i} ^ {\prime}\right) \leqslant \sum_ {i} \left(f \left(\mathrm{K} _ {i}\right) - f \left(\mathrm{K} _ {i} ^ {\prime}\right)\right).
\]

Let \((\mathbf{A}_i), (\mathbf{B}_i)\) be two finite families of subsets of \(\mathbf{X}\) , with the same index set, such that \(\mathbf{B}_i \subset \mathbf{A}_i\) for each \(i\) , \(f^* \left( \bigcup_i \mathbf{A}_i \right) < +\infty\) and \(f^*(\mathbf{B}_i) > -\infty\) for each \(i\) . Show that if \(f\) satisfies (AL), we have

\[
f ^ {*} \left(\bigcup_ {i} \mathrm{A} _ {i}\right) - f ^ {*} \left(\bigcup_ {i} \mathrm{B} _ {i}\right) \leqslant \sum_ {i} \left(f ^ {*} (\mathrm{A} _ {i}) - f ^ {*} (\mathrm{B} _ {i})\right).
\]

{[}Reduce to the case of an index set with two elements. Consider first the case where \(A_{i}\) and \(B_{i}\) are open sets, and use the following lemma: if \(U_{1}, U_{2}\) are two open sets and K is a compact set contained in \(U_{1} \cup U_{2}\) , then there exist compact sets \(K_{i} \subset U_{i}\) (i = 1, 2) such that \(K \subset K_{1} \cup K_{2}\) .{]} c) Suppose that f satisfies condition (AL). Show that for every increasing sequence \((\mathbf{A}_{n})\) of subsets of X such that \(f^{*}(\mathbf{A}_{n}) > -\infty\) , we have \(f^{*}\left(\bigcup_{n} \mathbf{A}_{n}\right) = \sup_{n} f^{*}(\mathbf{A}_{n})\) {[}use b){]}. Deduce that if f does not take the value \(-\infty\) on \(\Re(\mathbf{X})\) , then \(f^{*}\) is a capacity on X.

\begin{enumerate}
\def\labelenumi{\alph{enumi})}
\setcounter{enumi}{3}
\tightlist
\item
  Suppose that \(f\) satisfies the condition (AL). Show that the union of any sequence \((\mathbf{A}_n)\) of admissible sets such that \(f^*(\mathbf{A}_n) > -\infty\) for each \(n\) is admissible. {[}Use b) and c).{]}
\end{enumerate}

* 13) Let K be a compact subset of \(\mathbf{R}^2\) . For each real number \(y\) , let \(\delta_1(\mathbf{K}, y)\) {[}resp. \(\delta_2(\mathbf{K}, y)\) {]} denote the diameter of \(\mathbf{K} \cap ([0, +\infty [\times \{y\})]\) {[}resp. \(\mathbf{K} \cap (]-\infty, 0] \times \{y\})\) {]}. Show that the functions \(y \to \delta_1(\mathbf{K}, y)\) and \(y \to \delta_2(\mathbf{K}, y)\) are upper semi-continuous. Let \(\varphi\) be an increasing continuous mapping of \([0, +\infty]\) into itself such that \(\varphi(0) = 1\) and \(\varphi(+\infty) = 2\) , and let

\[
\psi (\mathbf {K}, y) = \varphi (\delta_ {1} (\mathbf {K}, y) \delta_ {2} (\mathbf {K}, y)) \quad \text { and } \quad f (\mathbf {K}) = \int_ {\mathrm{pr} _ {2} (\mathbf {K})} \psi (\mathbf {K}, y) d y.
\]

Show that \(f\) satisfies the conditions \((\mathbf{CA}_{\mathbf{I}})\) and \((\mathbf{CA}_{\mathbf{III}})\) and that we have \(f(\mathbf{K} \cup \mathbf{K}') \leqslant f(\mathbf{K}) + f(\mathbf{K}')\) for each pair of compact subsets \(\mathbf{K}, \mathbf{K}'\) of \(\mathbf{R}^2\) . But if \(\mathbf{A}\) is the closed set defined by \(x \geqslant 0\) and \(0 \leqslant y \leqslant 1\) , show that \(f_*(\mathbf{A}) = 1\) and \(f^*(\mathbf{A}) = 2\) .

\begin{enumerate}
\def\labelenumi{\arabic{enumi})}
\setcounter{enumi}{13}
\tightlist
\item
  Let \(\mathbf{X}\) be a metrizable space and let \(f\) be a capacity on \(\mathbf{X}\) such that \(f(\mathbf{A} \cup \mathbf{B}) \leqslant f(\mathbf{A}) + f(\mathbf{B})\) .
\end{enumerate}

\begin{enumerate}
\def\labelenumi{\alph{enumi})}
\item
  With the notation of Exercise 12, show that \(f^{*}(\mathbf{A} \cup \mathbf{B}) \leqslant f^{*}(\mathbf{A}) + f^{*}(\mathbf{B})\) and \(f_{*}(\mathbf{A} \cup \mathbf{B}) \leqslant f_{*}(\mathbf{A}) + f^{*}(\mathbf{B})\) for any two subsets \(\mathbf{A}, \mathbf{B}\) of \(\mathbf{X}\) . {[}If \(\mathbf{K}\) is a compact set contained in \(\mathbf{A} \cup \mathbf{B}\) , and \(\mathbf{U}\) is an open set containing \(\mathbf{B}\) , write \(\mathbf{K}\) in the form \(\mathbf{K} = (\mathbf{K} \cap \mathbf{U}) \cup (\mathbf{K} \cap \mathbf{C}\mathbf{U})\) .{]}
\item
  Let K be a compact subset of X having the power of the continuum, and let (A, B) be a partition of K into two sets such that every compact subset of A or of B is countable {[}cf.~§ 5, Exercise 18 d){]}. Show that if \(f(\{x\}) = 0\) for all \(x \in \mathbf{X}\) and if \(f(\mathbf{K}) > 0\) , then neither \(\mathbf{A}\) nor \(\mathbf{B}\) is capacitable with respect to \(f\) .
\end{enumerate}

* 15) Let \(\mu\) be the Lebesgue measure on \(\mathbf{R}\) ; define a capacity \(f\) on \(\mathbf{R}^2\) satisfying the condition (AL) of Exercise 12 by setting \(f(\mathbf{A}) = \mu^{*}(\mathrm{pr}_{1}\mathbf{A})\) (Proposition 15).

\begin{enumerate}
\def\labelenumi{\alph{enumi})}
\item
  Let A be a non-capacitable bounded subset of \(R^{2}\) (Exercise 14), let \(B_{0}\) be a circle \(\|x\| = r\) such that A is contained in the open disc \(\|x\| < r\) , and let \(B_{1}\) be a circle \(\|x\| = r'\) of radius \(r' > r\) . Show that \(A \cup B_{0}\) and \(A \cup B_{1}\) are capacitable, although their intersection A is not.
\item
  For each \(n\) , let \(\mathbf{C}_n\) be the set of all \(x \in \mathbb{R}^2\) such that
\end{enumerate}

\[
r <   | | \boldsymbol {x} | | <   r + \frac {\mathrm{I}}{n}.
\]

Show that each of the sets \(\mathbf{A} \cup \mathbf{C}_n\) is capacitable, although their intersection is not, and that we have \(\inf_n f(\mathbf{C}_n) \neq f\left(\bigcap_n \mathbf{C}_n\right)\) . *

¶ 16) Let X, X' be two metrizable spaces. A mapping \(f: \mathbf{X} \to \mathbf{X}'\) is said to be a Borel mapping if, for each closed subset F' of X', \(f^{-1}(\mathbf{F}')\) is a Borel set in X. For each even (resp. odd) countable ordinal \(\alpha\) , \(f\) is said to be of class \(\alpha\) if, for each closed subset F' of X', \(f^{-1}(\mathbf{F}')\) belongs to \(\mathfrak{S}_{\alpha}(\mathbf{X})\) {[}resp. \(\mathfrak{G}_{\alpha}(\mathbf{X})\) {]} (Exercise 4); then, for each open set G' in X', \(f^{-1}(\mathbf{G}')\) belongs to \(\mathfrak{G}_{\alpha}(\mathbf{X})\) {[}resp. \(\mathfrak{F}_{\alpha}(\mathbf{X})\) {]}. The Borel mappings of class 0 are precisely the continuous mappings.

\begin{enumerate}
\def\labelenumi{\alph{enumi})}
\item
  The characteristic function \(\varphi_{\mathbf{A}}\) of a subset \(\mathbf{A}\) of \(\mathbf{X}\) is of class \(\alpha\) if and only if \(\mathbf{A} \in \mathfrak{F}_{\alpha} \cap \mathfrak{G}_{\alpha}\) .
\item
  Let \(f: \mathbf{X} \to \mathbf{X}'\) be a mapping of class \(\alpha\) . Show that for every subset \(\mathbf{B}' \in \mathfrak{F}_{\beta}(\mathbf{X}')\) {[}resp. \(\mathbf{B}' \in \mathfrak{G}_{\beta}(\mathbf{X}')\) {]}, \(\overline{f}(\mathbf{B}')\) belongs to \(\mathfrak{F}_{\alpha + \beta}(\mathbf{X})\) {[}resp. \(\mathfrak{G}_{\alpha + \beta}(\mathbf{X})\) {]}. (Use transfinite induction on \(\beta\) .)
\item
  Let \(f: \mathbf{X} \to \mathbf{X}'\) be a mapping of class \(\alpha\) and let \(g: \mathbf{X}' \to \mathbf{X}''\) be a mapping of class \(\beta\) , where \(\mathbf{X}''\) is a third metrizable space. Show that \(g \circ f\) is a Borel mapping of class \(\alpha + \beta\) .
\item
  For each even (resp. odd) countable ordinal \(\alpha\) , let \((\mathbf{A}_n)\) be a sequence of subsets of \(\mathbf{X}\) which belong to \(\mathfrak{G}_{\alpha}\) (resp. \(\mathfrak{F}_{\alpha}\) ) and which cover \(\mathbf{X}\) . If a mapping \(f: \mathbf{X} \to \mathbf{X}'\) is such that \(f|_{\mathbf{A}_n}\) is of class \(\alpha\) for each \(n\) , then \(f\) is of class \(\alpha\) . Give the analogous result for a finite sequence of subsets of \(\mathbf{X}\) which belong to \(\mathfrak{F}_{\alpha}\) (resp. \(\mathfrak{G}_{\alpha}\) ) and cover \(\mathbf{X}\) .
\end{enumerate}

\begin{enumerate}
\def\labelenumi{\arabic{enumi})}
\setcounter{enumi}{16}
\tightlist
\item
  \begin{enumerate}
  \def\labelenumii{\alph{enumii})}
  \tightlist
  \item
    Let \(\mathbf{X}, \mathbf{X}'\) be two metrizable spaces, \(\mathbf{X}'\) being of countable type, and let \((\mathbf{U}_n')\) be a countable base for the topology of \(\mathbf{X}'\) . Let \(\alpha\) be an even (resp. odd) countable ordinal. If a mapping \(f: \mathbf{X} \to \mathbf{X}'\) is such that \(f^{-1}(\mathbf{U}_n')\) belongs to \(\mathfrak{G}_{\alpha}(\mathbf{X})\) {[}resp. \(\mathfrak{F}_{\alpha}(\mathbf{X})\) {]} for each \(n\) , then \(f\) is a Borel mapping of class \(\alpha\) .
  \end{enumerate}
\end{enumerate}

\begin{enumerate}
\def\labelenumi{\alph{enumi})}
\setcounter{enumi}{1}
\item
  Let \((\mathbf{X}_{n}^{\prime})\) be a sequence of metrizable spaces of countable type. A mapping \(f = (f_{n}) : \mathbf{X} \to \prod_{n} \mathbf{X}_{n}^{\prime}\) is of class \(\alpha\) if and only if each \(f_{n}\) is of class \(\alpha\) {[}use a){]}. Deduce that, if X is of countable type, the finite real-valued functions of class \(\alpha\) on X form a ring.
\item
  Let X, X' be two metrizable spaces of countable type, and let \(\alpha\) be an even (resp. odd) countable ordinal. If a mapping \(f: X \to X'\) is of class \(\alpha\) , show that its graph belongs to \(\mathfrak{F}_{\alpha}(X \times X')\) {[}resp. \(\mathfrak{G}_{\alpha}(X \times X')\) {]} {[}use b) and Exercise 16 b){]}.
\end{enumerate}

\begin{enumerate}
\def\labelenumi{\arabic{enumi})}
\setcounter{enumi}{17}
\tightlist
\item
  Let X be a Souslin space, \(X'\) a metrizable space of countable type and let f be a mapping of X into \(X'\) .
\end{enumerate}

\begin{enumerate}
\def\labelenumi{\alph{enumi})}
\item
  Show that if the graph G of f is a Souslin set in \(X \times X'\) , then f is a Borel function and therefore {[}Exercise 17 c){]} G is a Borel set (use Theorem 2, Corollary).
\item
  Show that if \(f\) is an injective Borel mapping and if \(\mathbf{X}'\) is a Souslin space, then the image under \(f\) of any Souslin set in \(\mathbf{X}\) is a Souslin set in \(\mathbf{X}'\) (use Exercise 17c)). If \(f\) is bijective, the inverse mapping is also a Borel mapping.
\end{enumerate}

\begin{enumerate}
\def\labelenumi{\arabic{enumi})}
\setcounter{enumi}{18}
\tightlist
\item
  Let X, X' be two metrizable spaces and let f be an almost open mapping (§ 5, Exercise 24) of X into X'. Show that if B' is any Borel set in X', then \(f^{-1}(B')\) is almost open in X (cf.~§ 5, Exercise 6). Deduce that if g is a Borel mapping of X' into a metrizable space X'\,', then g∘f is almost open. Give an example of a continuous mapping f and an almost open mapping g such that g∘f is not almost open {[}cf.~Exercise 10 b){]}.
\end{enumerate}

